\section{A matrix path with prescribed endpoints}\label{sec:path}

The action argument needs a symmetric matrix path from $K$ to $-K$.
Its differential equation also contains a lower triangular coefficient
matrix $C$, which we are free to choose.  We choose $C$ by matching the
terminal endpoint.  At the start of a deformation, the matching problem
has exactly one solution after identifying each row of $C$ with its
negative.  A compactness argument shows that this solution count cannot
disappear modulo two.

Fix $\delta>0$ throughout the remaining construction.  We use the
matrix field $M=M_\delta$ and scalar function $\mu=\mu_\delta$
constructed above.  Recall the properties needed here.  The field
$M$ is locally Lipschitz on $\Sym_N$, the space
of real symmetric $N\times N$ matrices.  The function $\mu$ is locally
Lipschitz and strictly increasing, with $\mu(0)=0$.  For
$K=\operatorname{diag}(k_1,\ldots,k_N)$ and every diagonal sign matrix
$S$,
\begin{equation}\label{path:equivariance}
 M(SBS)=SM(B)S.
\end{equation}
For every real unit vector $e$,
\begin{equation}\label{path:comparison}
 e^{\mathsf T}M(B)e
 \leq \mu\bigl(e^{\mathsf T}K^2e-(e^{\mathsf T}Be)^2\bigr).
\end{equation}
Finally, writing $e_i$ for the $i$th standard basis vector,
\begin{equation}\label{path:coordinate}
 Be_i=t e_i
 \quad\Longrightarrow\quad
 M(B)e_i=\mu(k_i^2-t^2)e_i.
\end{equation}
The equivariance in \eqref{path:equivariance} only concerns sign
matrices, which commute with the fixed matrix $K$.

\begin{proposition}[Connecting path]\label{path:existence}
Let $I=(x_-,x_+)$ be a bounded interval, let
$u\in H^1_0(I;\mathbb R^N)$ satisfy
\[
 \int_I u(x)u(x)^{\mathsf T}\,dx=I_N,
\]
and let $K=\operatorname{diag}(k_1,\ldots,k_N)$ with every $k_i>0$.
Suppose $M$ and $\mu$ have the regularity and properties
\eqref{path:equivariance}--\eqref{path:coordinate} just stated.
For every $\varepsilon>0$, there are a real lower triangular matrix
$C$ and a path $B\in C^1(\overline I;\Sym_N)$ such that,
with $a=Cu$,
\begin{equation}\label{path:equation}
 \varepsilon B'=M(B)-aa^{\mathsf T},
 \qquad B(x_-)=K,\qquad B(x_+)=-K.
\end{equation}
These choices satisfy
\begin{equation}\label{path:operator-bound}
 \|B(x)\|_{\mathrm{op}}\leq\|K\|_{\mathrm{op}}
 \qquad(x\in I).
\end{equation}
For fixed $M$, $K$, and $I$, the matrices $C$ can be chosen bounded
uniformly for $0<\varepsilon\leq1$.
\end{proposition}

The parity principle is the compact-one-manifold argument underlying
mod-two degree; see \cite[Section~4]{Milnor1965}. The smooth regular-value
facts used below are given in \cite[Sections~2--3]{Milnor1965}.
We prove the needed relative version for continuous families of bundle
sections here. A section of a vector bundle assigns to each base point a
vector in the fiber above it.  A zero is \emph{transverse} if the
derivative of the section in any local trivialization maps onto that
fiber.

\begin{lemma}[Continuation modulo two]\label{path:parity}
Let $Q$ be a smooth $d$-dimensional manifold without boundary, and let
$E\to Q$ be a smooth real vector bundle of rank $d$.  Suppose
$s_\theta$, $0\leq\theta\leq1$, is a continuous homotopy of sections and
\[
 \{(q,\theta)\in Q\times[0,1]:s_\theta(q)=0\}
\]
is compact.  If $s_0$ is smooth and has an odd number of zeros, all
transverse, then $s_1$ has a zero.
\end{lemma}

\begin{proof}
Suppose that $s_1$ has no zero.  Choose a relatively compact smooth
domain $\Omega\subset Q$ containing the projection of the entire zero
set.  Such a domain is obtained by taking a regular level of a smooth
compactly supported function equal to one near that projection.
Fix a smooth fiber metric on $E$.  On the compact set
\[
 (\partial\Omega\times[0,1])
 \ \cup\ (\overline\Omega\times\{1\}),
\]
the norm of the section has a strictly positive minimum.

Reparametrize the homotopy so that it is identically $s_0$ on a small
initial collar.  It can then be uniformly approximated on
$\overline\Omega\times[0,1]$ by a smooth section that agrees with $s_0$
on a smaller initial collar.  To construct this approximation, use
finitely many bundle charts, approximate the continuous coefficient
functions by smooth functions, and combine them with a smooth
partition of unity.  A cutoff supported inside the original collar
splices this approximation with the smooth section $s_0$.  Choose the
approximation error smaller than one quarter of the positive boundary
minimum.  The approximating section therefore has no zeros near the
spatial boundary or the final boundary.

We can also make this smooth section transverse to zero while keeping
it unchanged near the initial boundary.  Here are the finite-dimensional
details.  In an initial collar its zeros are already transverse, since
the derivative in the $Q$ directions is onto at each zero of $s_0$.
The remaining possible zeros lie in a compact subset of the interior.
Choose finitely many compactly supported local frame sections whose
values span every fiber on a neighborhood of this subset, with supports
away from the initial boundary and the forbidden boundaries.  Adding
linear combinations of these sections gives a finite-dimensional
family whose derivative in the new parameters is onto wherever
transversality is still needed. Restrict the parameters to a small open
ball in $\mathbb R^m$, where $m$ is the number of frame sections.
Smallness preserves the positive gap on the compact complement of the
spanning neighborhood and the initial collar, so no new zeros appear there.
The universal zero set in the interior is a smooth manifold of
dimension $m+1$. Sard's theorem applied to its projection onto
$\mathbb R^m$ provides arbitrarily small regular parameters. At a zero,
write the universal derivative as $Av+Db$, with $v$ a $(q,\theta)$ variation and
$b$ a parameter variation. Surjectivity of the projection on the kernel
means that $Av=-Db$ is solvable for every $b$. Where perturbations are
needed, $D$ is onto, so this is equivalent to $A$ being onto; elsewhere
the section is already transverse. Thus regular parameters make the
fixed-parameter section transverse.
Smallness in the smooth topology preserves the transversality already
present on the overlap with the initial collar, and smallness in norm
preserves the boundary gap.

The zero set of the resulting section is a compact smooth
one-dimensional manifold.  Its boundary consists exactly of the zeros
of $s_0$ on $\Omega\times\{0\}$.  A compact one-dimensional manifold
has an even number of boundary points, contradicting the assumed odd
count.  No orientation of $Q$ or $E$ is required.
\end{proof}



\begin{proof}[Proof of Proposition~\ref{path:existence}]
Put $\kappa=\|K\|_{\mathrm{op}}$ and $\ell=x_+-x_-$, and use the
continuous representative of $u$.  Choose a smooth nonnegative
function $\chi$ on $\Sym_N$, invariant under
orthogonal conjugation, compactly supported, and equal to one whenever
$\|B\|_{\mathrm{op}}\leq\kappa$.  For example, take
$\chi(B)=\eta(\operatorname{tr}B^2)$ with $\eta$ equal to one on
$[0,N\kappa^2]$ and zero outside a larger bounded interval.  The field
$\widehat M=\chi M$ is bounded and globally Lipschitz.

\smallskip
\noindent\emph{The terminal matching problem.}
Let $\mathcal L$ denote the vector space of real lower triangular
$N\times N$ matrices.  For $C\in\mathcal L$ and $0\leq\theta\leq1$,
solve
\begin{equation}\label{path:homotopy-ode}
 \varepsilon B'=\theta\widehat M(B)-(Cu)(Cu)^{\mathsf T},
 \qquad B(x_-)=K.
\end{equation}
The bounded globally Lipschitz field and continuous forcing give a
unique solution throughout $I$.  The integral equation and Gronwall's
inequality give continuous dependence on $(C,\theta)$.  Thus
\[
 \Phi_\theta(C)=B(x_+)+K
\]
is a continuous homotopy of maps from $\mathcal L$ to
$\Sym_N$.

At a zero of $\Phi_\theta$, integration of
\eqref{path:homotopy-ode} and orthonormality of $u$ give
\begin{equation}\label{path:coefficient-identity}
 CC^{\mathsf T}
 =2\varepsilon K+\theta\int_I\widehat M(B(x))\,dx.
\end{equation}
Consequently,
\begin{equation}\label{path:homotopy-bound}
 \|C\|_{\mathrm{HS}}^2
 \leq 2\varepsilon\operatorname{tr}K
       +\ell\sup_B|\operatorname{tr}\widehat M(B)|.
\end{equation}
For fixed $\varepsilon$, the full zero set in
$\mathcal L\times[0,1]$ is therefore compact.

\smallskip
\noindent\emph{Rows cannot vanish at a matching solution.}
Let $G$ be the group of diagonal sign matrices.  Equation
\eqref{path:equivariance}, invariance of $\chi$, and uniqueness of the
flow show that
\begin{equation}\label{path:terminal-equivariance}
 \Phi_\theta(SC)=S\Phi_\theta(C)S
 \qquad(S\in G).
\end{equation}
Suppose that row $i$ of $C$ vanishes, and let $S_i$ change only its
sign.  Then $S_iC=C$, so uniqueness gives $S_iBS_i=B$ throughout the
flow.  Hence $Be_i=b_i e_i$.  By \eqref{path:coordinate},
\[
 \varepsilon b_i'
 =\theta\chi(B(x))\mu(k_i^2-b_i^2),
 \qquad b_i(x_-)=k_i.
\]
Regarding $\theta\chi(B(x))/\varepsilon$ as a known continuous
coefficient, local Lipschitz uniqueness for this scalar equation gives
$b_i\equiv k_i$.  Since $k_i>0$, the terminal value cannot be $-k_i$.
Thus no zero of $\Phi_\theta$ has a vanishing row.

Let $\mathcal L^\times\subset\mathcal L$ be the open subset on which
every row is nonzero.  The compact joint zero set lies inside
$\mathcal L^\times\times[0,1]$, and hence a positive distance away
from every zero-row stratum.  We do not require $C$ to be invertible.

\smallskip
\noindent\emph{One zero modulo signs persists.}
The action of $G$ on $\mathcal L^\times$ is free: changing the sign
of a nonzero row cannot fix $C$.  Its quotient
\[
 Q=\mathcal L^\times/G
\]
is therefore a smooth manifold of dimension $d=N(N+1)/2$.  Concretely,
each row $i$ belongs to $\mathbb R^i\setminus\{0\}$ and is identified
with its negative.  Sign conjugation on symmetric matrices defines
the rank-$d$ vector bundle
\[
 E=(\mathcal L^\times\times\Sym_N)/G
 \longrightarrow Q,
 \qquad
 (C,Z)\sim(SC,SZS).
\]
By \eqref{path:terminal-equivariance}, the maps $\Phi_\theta$ descend
to a continuous homotopy of sections of this bundle.  Their joint
zero set remains compact.

At $\theta=0$,
\[
 \Phi_0(C)=2K-\varepsilon^{-1}CC^{\mathsf T}.
\]
The lower triangular solutions of $CC^{\mathsf T}=2\varepsilon K$
are exactly
\[
 C=\operatorname{diag}
   (\pm\sqrt{2\varepsilon k_1},\ldots,
    \pm\sqrt{2\varepsilon k_N}).
\]
Indeed, the first row determines $C_{11}$, the other entries of the
first column then vanish, and the same argument applies successively
to the remaining principal blocks.  These solutions form one $G$-orbit.
At any one of them, writing $C=\operatorname{diag}(c_i)$,
\[
 D\Phi_0(C)[H]
 =-\varepsilon^{-1}(HC^{\mathsf T}+CH^{\mathsf T})
 \qquad(H\in\mathcal L).
\]
Its diagonal entries are $-2c_iH_{ii}/\varepsilon$, and its entries
with $i>j$ are $-c_jH_{ij}/\varepsilon$.  Since every $c_i\ne0$, this
linear map is an isomorphism onto the symmetric matrices.  The
initial quotient section has exactly one transverse zero.
Lemma~\ref{path:parity} now gives a zero of $\Phi_1$.

\smallskip
\noindent\emph{The endpoints remove the cutoff.}
Take a path obtained from this zero.  For a fixed real unit vector
$e$, put $h(x)=e^{\mathsf T}B(x)e$.  Whenever $|h(x)|>\kappa$,
\eqref{path:comparison} implies
\begin{align*}
 \varepsilon h'(x)
 &\leq\chi(B(x))
    \mu\bigl(e^{\mathsf T}K^2e-h(x)^2\bigr)
       -|e^{\mathsf T}Cu(x)|^2
 \leq0.
\end{align*}
Here $e^{\mathsf T}K^2e\leq\kappa^2<h(x)^2$; the last inequality
uses $\chi\geq0$ and remains true where $\chi=0$.
Both endpoint values of $h$ lie in $[-\kappa,\kappa]$.
A component of $\{h>\kappa\}$ cannot start from the initial side:
on that component $h$ is nonincreasing, whereas its left endpoint
has value $\kappa$.  Similarly, a component of $\{h<-\kappa\}$
cannot return to the final endpoint: its right endpoint would have
value $-\kappa$, again contradicting nonincrease.  Thus
$|h(x)|\leq\kappa$ throughout $I$.

Since $e$ was arbitrary, this proves
\eqref{path:operator-bound}.  Therefore $\chi(B(x))=1$ everywhere,
and the path solves the uncut equation \eqref{path:equation}.
Finally, its integrated equation gives, for $0<\varepsilon\leq1$,
\begin{equation}\label{path:uniform-coefficients}
 \|C\|_{\mathrm{HS}}^2
 \leq 2\operatorname{tr}K
 +\ell\sup_{\|B\|_{\mathrm{op}}\leq\kappa}
             |\operatorname{tr}M(B)|.
\end{equation}
This is the asserted bound on $C$.
\end{proof}

For the regularized field $M_\delta$, the bound
\eqref{path:uniform-coefficients} may depend on $\delta$.
With $\delta$ fixed, it implies that all pairs $(B(x),Cu(x))$ produced
for $0<\varepsilon\leq1$ lie in one compact subset of
$\Sym_N\times\mathbb R^N$.
This is the compactness needed in the action limit; it requires no
bound on $B'$ uniform in $\varepsilon$.
